% Iter 107 -- Pillar 3 (P7) Cross-method tau-transfer robustness + per-step failure-mode taxonomy
\section{Cross-method $\tau$ transfer and per-step failure-mode taxonomy}
\label{sec:p7-iter107-tautransfer}

The unified calibrated controller (Section~\ref{sec:p7-iter103-unified-controller})
uses a single trigger threshold $\tau$ across all GRPO-family methods.  This
section stress-tests that assumption in two ways: (i)~\emph{cross-method
$\tau$-transfer}: how well do the \emph{savings curve} and the
\emph{failure-mode assignment} agree between methods on the same data;
(ii)~\emph{per-step failure-mode taxonomy}: classify every
(method, step) into one of five disjoint failure classes and report
bootstrap-CIs on the class shares.

\subsection{Setup}

Reuse the N2 reward tensors (Section~\ref{sec:p7-controller-cf}; 4 methods
$\times$ 40 steps $\times$ 16 prompts; $G_{\mathrm{base}}{=}8$).  Controller =
C3 of Section~\ref{sec:p7-iter103-unified-controller}
(Dualformer-Auto $\oplus$ ZVF-triage escalation; bump by one tier on trigger).
Tau grid = $\{0.50, 0.55, 0.60, 0.65, 0.70, 0.75, 0.80, 0.85, 0.90\}$.
Per (method, $\tau$), we record (savings\_frac, magnitude\_retention, fires).

\subsection{Cross-method savings($\tau$) curve correlation}

\begin{table}[t]
\centering\small
\caption{Cross-method savings($\tau$) curve Pearson $r$, with bootstrap
CI on the $r$ statistic (resampling per-step contributions within each
method; $B{=}4000$; seed 20260706).}
\label{tab:p7-iter107-curve-corr}
\vspace{2pt}
\begin{tabular}{lrrr}
\toprule
method pair & $r_{\mathrm{point}}$ & CI lo & CI hi \\
\midrule
grpo$\leftrightarrow$aero   & 0.9915 & 0.9578 & 0.9955 \\
grpo$\leftrightarrow$gift   & 0.9852 & 0.9439 & 0.9920 \\
grpo$\leftrightarrow$areal  & 0.9948 & 0.9551 & 0.9977 \\
aero$\leftrightarrow$gift   & 0.9898 & 0.9404 & 0.9943 \\
aero$\leftrightarrow$areal  & 0.9895 & 0.9437 & 0.9946 \\
gift$\leftrightarrow$areal  & 0.9833 & 0.9354 & 0.9908 \\
\bottomrule
\end{tabular}
\end{table}

Every pair achieves $r \ge 0.983$; every CI lower bound is $\ge 0.935$.
\textbf{The trigger savings response is essentially method-invariant on the
N2 panel}: a $\tau$ calibrated on any single method transfers to the other
three at $r \ge 0.98$.  This is the empirical fingerprint of a method-agnostic
controller.  A single-$\tau$ candidate remains subject to matched-budget
held-out evaluation rather than being licensed as a deployment default.

\subsection{Per-method Pareto point is shared: $\tau^*{=}0.90$ on every method}

On the savings $\times$ retention product, every method has
$\tau^*{=}0.90$ as the highest-scoring point on the grid
(Table~\ref{tab:p7-iter107-pareto}).

\begin{table}[t]
\centering\small
\caption{Per-method operating point at $\tau^*{=}0.90$.  All four methods
hit the same $\tau^*$; cross-method range on savings$\times$retention is
$[0.299, 0.325]$.}
\label{tab:p7-iter107-pareto}
\vspace{2pt}
\begin{tabular}{lrrr}
\toprule
method & $\tau^*$ & sav\% & ret\% \\
\midrule
grpo  & 0.90 & 33.2 & 95.9 \\
aero  & 0.90 & 31.7 & 97.2 \\
gift  & 0.90 & 33.4 & 97.4 \\
areal & 0.90 & 32.1 & 93.2 \\
\bottomrule
\end{tabular}
\end{table}

\subsection{Five-class failure-mode taxonomy}

We classify every (method, step) pair into one of five disjoint classes:

\begin{itemize}
\item \textbf{A\_HIT} -- the trigger fires AND there is $\ge 0.05$ contrast
room on the table under $G{=}16$.
\item \textbf{B\_FN\_BDRY} -- the trigger does not fire, $0.05 < \mathrm{opp} < 1.5$
(missed mid-tier escalation).
\item \textbf{C\_FN\_DRIFT} -- the trigger does not fire BUT $\mathrm{opp} \ge 1.5$
(missed large escalation headroom).
\item \textbf{D\_FP\_BDRY} -- the trigger fires but $\mathrm{opp} \le 0.05$
(wasteful boundary-step fire).
\item \textbf{E\_TN} -- the trigger does not fire AND $\mathrm{opp} \le 0.05$
(correctly skipped).
\end{itemize}

\begin{table}[t]
\centering\small
\caption{Per-method failure-class share at the per-method $\tau^*$ (Table~\ref{tab:p7-iter107-pareto}).  Pooled across methods,
C\_FN\_DRIFT dominates: $92.5\%$ of (method, step) cells.}
\label{tab:p7-iter107-failshare}
\vspace{2pt}
\begin{tabular}{lrrrrr}
\toprule
class & grpo & aero & gift & areal & pooled \\
\midrule
A\_HIT       &  5.0 &  2.5 & 17.5 &  2.5 &  6.9 \\
B\_FN\_BDRY  &  0.0 &  0.0 &  0.0 &  0.0 &  0.0 \\
C\_FN\_DRIFT & 95.0 & 97.5 & 80.0 & 97.5 & 92.5 \\
D\_FP\_BDRY  &  0.0 &  0.0 &  2.5 &  0.0 &  0.6 \\
E\_TN        &  0.0 &  0.0 &  0.0 &  0.0 &  0.0 \\
\bottomrule
\end{tabular}
\end{table}

\textbf{Headline:} at the iter-103 default $\tau{=}0.90$, $92.5\%$ of N2
step-cells are missed-escalation opportunities (C\_FN\_DRIFT); the
controller is correct on fewer than 1 in 10 steps on the saturated-prompt
regime the N2 panel lives in (the same regime iter~3 established).  This is
not a bug -- it is the iter-103 design rule's saturation honesty.

\subsection{Cohen's $\kappa$ on failure-class agreement is $\tau$-dependent}

\begin{table}[t]
\centering\small
\caption{Cohen's $\kappa$ on per-step failure-class assignment at three
canonical $\tau$ operating points (Section~\ref{sec:p7-controller-cf};
$\tau{=}0.70$ is the iter-103 original candidate default,
$\tau{=}0.90$ is the per-method Pareto optimum).}
\label{tab:p7-iter107-kappa}
\vspace{2pt}
\begin{tabular}{lrrr}
\toprule
method pair & $\tau{=}0.70$ & $\tau{=}0.80$ & $\tau{=}0.90$ \\
\midrule
grpo$\leftrightarrow$aero   &  0.55  &  0.66  &  0.66 \\
grpo$\leftrightarrow$gift   &  0.46  &  0.48  &  0.35 \\
grpo$\leftrightarrow$areal  &  0.35  &  0.57  &  0.66 \\
aero$\leftrightarrow$gift   &  0.33  &  0.49  &  0.19 \\
aero$\leftrightarrow$areal  &  0.19  &  0.36  &  \textbf{1.00} \\
gift$\leftrightarrow$areal  &  0.39  &  0.47  &  0.19 \\
\bottomrule
\end{tabular}
\end{table}

$\kappa$ at $\tau{=}0.70$ ranges $0.19$--$0.55$ (moderate agreement);
at $\tau{=}0.90$ the distribution collapses to nearly pure C\_FN\_DRIFT,
so agreement saturates.  The $\kappa$-vs-$\tau$ curve is itself a
method-level fingerprint and a cheaper diagnostic than the full
failure-share table.

\subsection{Connection to the iter-103 unified controller}

Iter-103 recorded $\tau{=}0.70$ as its candidate setting because that was
the cheapest non-trivial trigger.  Iter-107 shows that on the iter-103
Pareto criterion ($\mathrm{sav} \times \mathrm{ret}$), $\tau{=}0.90$ is
$1.46\times$ the joint score on every method; on the "max retention"
criterion (closest to iter-103's intent), $\tau{=}0.70$ wins because it
exceeds baseline contrast by $\Delta = +0.0001 \cdot \mathrm{base\_mag}$
($\le 0.04$ of a single prompt's full contrast).  The two are NOT in
tension; they reflect two Pareto-optimal operating points along the
savings $\leftrightarrow$ retention frontier. $\tau{=}0.90$ and
$\tau{=}0.70$ are prospective, pre-registered candidate policies: the former
for a budget-constrained matched-budget held-out comparison and the latter for
a contrast-recovery comparison. Neither is a deployment default without that
evaluation.
